Transport inventories are usually built bottom-up from a vehicle fleet and reconciled
against fuel supply afterwards, resting their weakest input---fleet activity---under their
most-used output. We invert that for Pakistan: road transport (IPCC 1A3b) and mixed
off-road machinery are anchored on the road customer class of OCAC's sectoral petroleum
sales for FY2023--24 (\SI{7209}{\kilo\tonne} petrol, \SI{5551}{\kilo\tonne} diesel; those
customer classes sum to the national total exactly, excluding bulk industrial, utility,
government and railway diesel by construction) and on metered transport CNG. The fleet
model distributes this fuel and cannot size it: a bottom-up build on the same data would
miss diesel by half and overstate CNG threefold. Class shares are vehicle-kilometres
weighted by declared energy intensities and normalised within each fuel, so that judgement
moves composition and never magnitude---correcting it from distance-only weighting moved
two-wheelers from \SI{42.3}{} to \SI{34.1}{\percent} of on-road \COtwo{}
(\SI{13.40}{\mega\tonne}) while the sector total
moved by \SI{0.05}{\mega\tonne}. We report \SI{40.19}{\mega\tonne}~\COe{} as the
\emph{conditional road residual under the current claimant assignment} and
\SI{1.73}{\mega\tonne} for mixed off-road machinery, with a two-tier ledger that closes
both controls exactly and carries the unresolved retail-channel claimants as status rows,
priced at 2.5--10\% of the road row and treated as transfers to other IPCC categories
rather than reductions. An independent railway cross-check agrees to \SI{3.8}{\percent},
and against a fiscal-year-blended EDGAR benchmark the estimate is \num{0.85} of
it---a scope-approximate diagnostic, since EDGAR's transport boundary is inferred
rather than published. The
reported interval is $\pm\SI{11.4}{\percent}$, conditional on the boundary; a ten-year
fuel-based series and per-fuel provincial allocations follow at no further modelling cost.
